\documentclass[11pt]{amsart}
\usepackage[T1]{fontenc}
\usepackage{amsmath,amssymb,amsthm}
\usepackage{graphicx}
\usepackage[colorlinks=true,linkcolor=blue,citecolor=blue,urlcolor=blue]{hyperref}

\newtheorem{theorem}{Theorem}
\newtheorem{lemma}{Lemma}
\newcommand{\T}{\mathbb T}
\newcommand{\muN}{\mu_N}

\begin{document}

\title{The $L^1$-Gradient Endpoint for Badly Approximable Kronecker Sequences}
\author{Run fa\_banach\_001}
\date{August 11, 2026}
\maketitle
\pagestyle{plain}

\section*{Status}

This is a candidate full answer, subject to human review, to both questions
following Theorem 7 of Brown and Steinerberger \cite{BS}: the regular-grid
$L^1$-gradient estimate does hold uniformly for every badly approximable
Kronecker sequence, so such an infinite sequence exists.

\section*{Source Question}

For $d\geq2$ and a badly approximable vector
$\alpha\in\mathbb R^d$, put $x_j=j\alpha\pmod1$ in $\T^d$.  The source
proves, uniformly in $N$, that
\[
 \left|\int_{\T^d}f-\frac1N\sum_{j=1}^N f(x_j)\right|
 \lesssim_{\alpha,d}N^{-1/d}
 \|\nabla f\|_\infty^{(d-1)/d}\|\nabla f\|_2^{1/d}.
\]
For a regular grid it proves the stronger endpoint with
$\|\nabla f\|_1$ in place of $\|\nabla f\|_2$, and asks whether the
same is true for the Kronecker sequence.  More generally, it asks whether
any single infinite sequence attains that endpoint uniformly in $N$.

\begin{center}
\includegraphics[width=\textwidth]{figures/open_problem_crop.png}
\end{center}

\section*{Proof Intuition}

The source splits the empirical measure into a short heat-flow part and a
smoothed part.  It bounds the first only by transportation distance, forcing
an $L^\infty$ gradient, and bounds the second in $\dot H^{-1}$, forcing an
$L^2$ gradient.  Keep the same split, but represent both pieces as
divergences of vector fields.

The low-frequency field has $L^2$ norm $O(N^{-1/d})$ by the source's own
Fourier estimate.  The high-frequency field is a sum of truncated Coulomb
kernels.  Bad approximability separates the first $N$ orbit points at scale
$N^{-1/d}$; a packing argument then puts this sum in the critical weak space
$L^{d/(d-1),\infty}$ with the same norm.  Lorentz H\"older duality against
$L^{d,1}$, followed by the elementary interpolation
\[
 \|g\|_{L^{d,1}}
 \lesssim_d \|g\|_\infty^{(d-1)/d}\|g\|_1^{1/d},
\]
gives exactly the missing endpoint.

\section*{Critical Heat-Flow Field}

Write $h=N^{-1/d}$ and
\[
 \muN=\frac1N\sum_{j=1}^N\delta_{x_j}.
\]
Let $P_t=e^{t\Delta}$ be heat flow on $\T^d$, with kernel $p_t$, and
take $t=h^2$.  We use Lorentz spaces with respect to normalized Lebesgue
measure.

\begin{lemma}[Separated heat atoms]\label{lem:high}
Suppose $x_1,\ldots,x_N\in\T^d$ satisfy
\[
 \min_{i\ne j}d(x_i,x_j)\geq a h.
\]
There is a vector field $F_{\rm hi}$ such that
\[
 \operatorname{div}F_{\rm hi}=\muN-P_t\muN
 \quad\text{and}\quad
 \|F_{\rm hi}\|_{L^{d/(d-1),\infty}}
 \leq C_{a,d}h.
\]
\end{lemma}

\begin{proof}
In distributions, set
\[
 F_{\rm hi}(x)
 =-\frac1N\sum_{j=1}^N K_t(x-x_j),
 \qquad
 K_t(z)=\int_0^t\nabla p_s(z)\,ds.
\]
Since $\partial_sp_s=\Delta p_s$, taking the divergence gives the asserted
identity.  The standard Gaussian gradient estimate for the torus heat kernel
implies, with $r=d(z,0)$,
\begin{equation}\label{eq:kernel}
 |K_t(z)|\leq C_d r^{1-d}
 \exp\!\left(-\frac{r^2}{C_dt}\right).
\end{equation}

Let $r(x)=\min_jd(x,x_j)$.  Divide the points into annuli of width $h$
around $x$.  Separation gives $O_{a,d}(m^{d-1})$ points in the $m$th
annulus.  Summing \eqref{eq:kernel}, and treating the nearest point
separately, yields
\begin{equation}\label{eq:pointwise}
 |F_{\rm hi}(x)|
 \leq C_{a,d}\bigl(h+h^d r(x)^{1-d}\bigr).
\end{equation}

Put $q=d/(d-1)$.  If $\lambda>2C_{a,d}h$, then
\eqref{eq:pointwise} shows that
\[
 \{|F_{\rm hi}|>\lambda\}
 \subseteq\bigcup_{j=1}^N B(x_j,R_\lambda),
 \qquad
 R_\lambda=C_{a,d}
 \left(\frac{h^d}{\lambda}\right)^{1/(d-1)}.
\]
Consequently
\[
 |\{|F_{\rm hi}|>\lambda\}|
 \leq C_d N R_\lambda^d
 \leq C_{a,d}\left(\frac h\lambda\right)^q.
\]
For $\lambda\leq2C_{a,d}h$ the trivial measure bound gives the same
weak-norm conclusion.  Hence
$\sup_{\lambda>0}\lambda
|\{|F_{\rm hi}|>\lambda\}|^{1/q}\lesssim_{a,d}h$.
\end{proof}

\begin{lemma}[The smoothed Kronecker field]\label{lem:low}
If $x_j=j\alpha\pmod1$ and $\alpha$ is badly approximable, there is a
vector field $F_{\rm lo}$ such that
\[
 \operatorname{div}F_{\rm lo}=P_t\muN-1
 \quad\text{and}\quad
 \|F_{\rm lo}\|_2\leq C_{\alpha,d}h.
\]
In particular,
$\|F_{\rm lo}\|_{L^{d/(d-1),\infty}}\leq C_{\alpha,d}h$.
\end{lemma}

\begin{proof}
Take $F_{\rm lo}=\nabla u$, where
$\Delta u=P_t\muN-1$ and $u$ has mean zero.  Thus
\[
 \|F_{\rm lo}\|_2
 =\|P_t\muN-1\|_{\dot H^{-1}}.
\]
The dyadic Fourier estimate proved in \cite[proofs of Theorems 5 and 6]{BS}
is
\[
 \|P_t\muN-1\|_{\dot H^{-1}}
 \leq \frac{C_{\alpha,d}}{N t^{(d-1)/2}}.
\]
At $t=h^2=N^{-2/d}$ the right side is $C_{\alpha,d}h$.
Finally $d/(d-1)\leq2$, so on the probability space $\T^d$ the $L^2$
bound implies the asserted weak-Lorentz bound.
\end{proof}

For completeness, the other input needed by Lemma \ref{lem:high} follows
directly from the source definition of bad approximability.  If
$1\leq i<j\leq N$, then
\[
 d(x_i,x_j)
 \geq\max_{1\leq\ell\leq d}
 \|(j-i)\alpha_\ell\|_{\mathbb R/\mathbb Z}
 \geq c_\alpha(j-i)^{-1/d}
 \geq c_\alpha h.
\]

\section*{Full Answer}

\begin{theorem}\label{thm:answer}
Let $d\geq2$ and let $\alpha\in\mathbb R^d$ be badly approximable.  For
$x_j=j\alpha\pmod1$ and every differentiable
$f:\T^d\to\mathbb R$,
\[
 \left|\int_{\T^d}f(x)\,dx-\frac1N\sum_{j=1}^Nf(x_j)\right|
 \leq C_{\alpha,d}N^{-1/d}
 \|\nabla f\|_\infty^{(d-1)/d}
 \|\nabla f\|_1^{1/d}.
\]
The constant is independent of $N$.  Hence both questions following
Theorem 7 of \cite{BS} have positive answers.
\end{theorem}

\begin{proof}
Let $F=F_{\rm hi}+F_{\rm lo}$ be supplied by Lemmas
\ref{lem:high} and \ref{lem:low}.  Then
\[
 \operatorname{div}F=\muN-1,
 \qquad
 \|F\|_{L^{d/(d-1),\infty}}\leq C_{\alpha,d}h.
\]
Distributional integration by parts and O'Neil's Lorentz H\"older inequality
\cite{ONeil} give
\begin{align*}
 \left|\int f\,d(\muN-dx)\right|
 &\leq C_d
 \|F\|_{L^{d/(d-1),\infty}}
 \|\nabla f\|_{L^{d,1}}\\
 &\leq C_{\alpha,d}h\|\nabla f\|_{L^{d,1}}.
\end{align*}
If $M=\|\nabla f\|_\infty$ and $L=\|\nabla f\|_1$, the layer-cake
formula and H\"older's inequality give
\begin{align*}
 \|\nabla f\|_{L^{d,1}}
 &=d\int_0^M
   |\{|\nabla f|>\lambda\}|^{1/d}\,d\lambda\\
 &\leq dM^{(d-1)/d}
 \left(\int_0^M|\{|\nabla f|>\lambda\}|\,d\lambda\right)^{1/d}\\
 &=dM^{(d-1)/d}L^{1/d}.
\end{align*}
Since $h=N^{-1/d}$, this proves the claim.
\end{proof}

\section*{Verification Notes}

The proof does not infer an $L^1$ endpoint from the source's $W_2$ estimate;
that would be invalid.  It retains the divergence fields before applying
duality.  The two pieces live in the same critical weak space for different
reasons: $L^2\subset L^{d/(d-1),\infty}$ for the smoothed part, while
separation controls the Coulomb singularities of the unsmoothed part.

The powers can be checked from the distribution calculation.  A single
weighted Coulomb singularity has size $N^{-1}r^{1-d}=h^dr^{1-d}$; $N$
such singularities produce distribution
$O((h/\lambda)^{d/(d-1)})$, whose weak norm is exactly $O(h)$.

Smooth functions suffice for the distributional computation.  Standard
mollification extends the estimate to differentiable/Lipschitz functions
with gradient in $L^1\cap L^\infty$, which is the source's class.

\section*{Novelty Check and Scope}

A bounded local and web search on August 11, 2026 used the arXiv identifier,
the exact endpoint sentence, and combinations of ``Kronecker sequence'',
``$L^1$ gradient'', ``Lorentz'', ``numerical integration'', and ``badly
approximable''.  It found the source and its journal version, later papers
on Wasserstein equidistribution, and 2024--2026 work on quasi-uniformity of
badly approximable Kronecker sequences.  It found no paper claiming this
endpoint or using the critical heat-field argument.  Novelty is plausible,
not certified.

The theorem answers the explicit $L^1$ question for the same
badly approximable vectors and dimensions $d\geq2$ treated by the source.
It does not characterize all sequences satisfying the endpoint, nor decide
whether bad approximability is necessary.

\section*{Human Review Recommendation}

Check the annular-packing estimate \eqref{eq:pointwise} and the sign in the
heat-flow divergence split.  Then compare Lemma \ref{lem:low} directly with
the Fourier sum in the source's proofs of Theorems 5 and 6.  If those pass,
Theorem \ref{thm:answer} solves both stated questions.

\begin{thebibliography}{9}

\bibitem{BS}
L. Brown and S. Steinerberger,
\emph{On the Wasserstein distance between classical sequences and the
Lebesgue measure},
Trans. Amer. Math. Soc. \textbf{373} (2020), 8943--8962;
arXiv:1909.09046.

\bibitem{ONeil}
R. O'Neil,
\emph{Convolution operators and $L(p,q)$ spaces},
Duke Math. J. \textbf{30} (1963), 129--142.

\end{thebibliography}

\end{document}
